\problemname{Range Sum 1}
Du har en lista med heltalen $A_0, A_1, \dots, A_{N-1}$.

Du ska svara på $Q$ stycken frågor ''beräkna $A_l+A_{l+1}+\cdots+A_{r}$'', där $l$ och $r$ är given
för varje fråga. 

\section*{Indata}
Den första raden av indata innehåller två heltal $N, Q$ ($1 \leq N, Q \leq 2 \cdot 10^5$).

Nästa rad innehåller $N$ positiva heltal $A_0, A_1, \dots, A_{N-1}$ ($1 \leq A_i \leq 10^9$).

Därefter följer $Q$ rader, som vardera innehåller två heltal $L$ och $R$ ($0 \leq L \leq R \leq N-1$),
intervallet för varje fråga.

\section*{Utdata}
Skriv ut $Q$ rader som vardera innehåller ett heltal, svaret på motsvarande fråga.

\section*{Poängsättning}
Din lösning kommer att testas på en mängd testfallsgrupper.
För att få poäng för en grupp så måste du klara alla testfall i gruppen.

\noindent
\begin{tabular}{| l | l | p{12cm} |}
  \hline
  \textbf{Grupp} & \textbf{Poäng} & \textbf{Gränser} \\ \hline
  $1$    & $50$      & $N, Q \leq 1000$ \\ \hline
  $2$    & $50$      & Inga ytterligare begränsningar. \\ \hline
\end{tabular}

\section*{Förklaring av exempelfall 1}
I första frågan är $L=0$ och $R=3$. Därmed är svaret $A_0+A_1+A_2+A_3=1+4+3+1=9$.
